\section{Phase 2: Geographic Voronoi Districts}
\label{sec:phase2}

\subsection{Motivation: When Graph Distance Diverges from Geography}

Phase~1 \cvd{} measures proximity via BFS hop count on the census-tract
adjacency graph.
BFS hop count is a faithful proxy for geographic distance in most of the
contiguous United States: inland states with roughly regular census-tract
shapes produce adjacency graphs where hop count is proportional to road
distance, and Voronoi regions based on hop count are reasonable geographic
approximations.

However, BFS hop count diverges from geographic Euclidean distance in two
important cases:

\paragraph{Case 1: Coastal and waterfront states.}
States with substantial water boundaries (Florida, Louisiana, Michigan,
Maryland) have census tracts that are geographically far apart but connected
via many intermediate tracts along land corridors.
Two tracts on opposite shores of a bay may be 2 miles apart in Euclidean
distance but 50 hops apart in the adjacency graph (via the land route).
BFS Voronoi in such states assigns coastal tracts to ``wrong'' seeds based
on graph topology rather than actual geographic proximity, producing
elongated peninsular districts rather than compact coastal ones.

\paragraph{Case 2: Elongated rural states.}
States with elongated geographic shapes and irregular census-tract sizes
(Nevada, Montana, Wyoming) have rural tracts that are large in area but
contribute only 1 hop to BFS distance.
A rural tract covering 1{,}000 square miles is treated the same as an urban
tract covering 1 square mile in hop-count space.
BFS Voronoi places district boundaries in locations that are geometrically
natural in hop-count space but geographically arbitrary.

Phase~2 \cvd{} addresses both cases by replacing BFS hop count with
Euclidean distance on projected tract centroids.

\subsection{Geographic CVD: Algorithm Modification}

The Phase~2 modification to Algorithm~\ref{alg:cvd} is minimal:

\begin{enumerate}
\item \textbf{Data requirement}: Load a \texttt{tract\_centroids} field in
  \texttt{LoadedGraph}: \texttt{pub tract\_centroids: Vec<(f64, f64)>},
  containing the (longitude, latitude) of each tract's polygon centroid.
  Centroids are available from TIGER/Line shapefiles as the centroid of
  each tract polygon, derivable from the \texttt{\_geoids.json} file
  already present in the pipeline.

\item \textbf{Metric change}: Replace $d_{\mathrm{BFS}}(t, \mathrm{seed}_j)$
  with $d_{\mathrm{Euc}}(t, \mathrm{seed}_j)$, the Euclidean distance between
  projected centroids.

\item \textbf{Centroid update}: Replace the graph-distance medoid with the
  population-weighted centroid:
  \[
    \mathrm{centre}_j = \frac{\sum_{t \in D_j} p_t \cdot \mathbf{x}_t}
                             {\sum_{t \in D_j} p_t},
  \]
  where $\mathbf{x}_t = (x_t, y_t)$ is the projected centroid of tract $t$.
  This is the true Euclidean centroid, not a medoid.
  The nearest tract to the centroid (by Euclidean distance) is used as the
  new seed for subsequent Voronoi assignments.

\item \textbf{Projection}: Use EPSG:5070 Albers Equal Area Conic for all
  50 states~\citep{epsg5070}. For continental US states, raw (lon, lat)
  Euclidean distance is adequate; for Alaska and Hawaii, where meridian
  convergence is severe, EPSG:5070 is mandatory to avoid distortion.
  Using EPSG:5070 uniformly ensures consistency with the projection used
  in \texttt{redist-map} for area and compactness computations.
\end{enumerate}

\subsection{Expected Improvement over Phase 1}

Phase~2 geographic \cvd{} is expected to yield a 5--10\% Polsby-Popper
improvement over Phase~1 graph-distance \cvd{} for states where BFS
hop count diverges substantially from geographic distance.

The intuition: Euclidean Voronoi regions on a geographic plane are exactly
the geometric cells that minimise mean squared distance from seed to tract
centroid. These cells are circular in uniform-density regions, and elliptical
in anisotropic regions, both of which are more compact (higher PP) than the
irregular BFS Voronoi regions.
For states where hop count faithfully proxies geography (Iowa, Kansas),
the Phase~2 improvement is expected to be small ($<3\%$).
For coastal states (Florida, Louisiana), the improvement is expected to be
larger (5--10\%) because Phase~2 avoids the along-shore bias of hop-count
Voronoi.

\subsection{Implementation Plan}

Phase~2 \cvd{} is not yet implemented; it is deferred pending the addition
of \texttt{tract\_centroids} to \texttt{LoadedGraph}.
The implementation plan is:

\begin{enumerate}
\item \textbf{Add \texttt{tract\_centroids} to \texttt{LoadedGraph}}:
  Derive centroid (lon, lat) for each tract from the existing TIGER/Line
  shapefile adjacency loader.
  Store as a \texttt{Vec<(f64, f64)>} field in the \texttt{LoadedGraph} struct.
  Add to the \texttt{.pkl} adjacency cache for zero-cost access on subsequent runs.

\item \textbf{Implement \texttt{VoronoiMetric::Geographic}}:
  Add the Euclidean distance branch to the \cvd{} Rust implementation.
  Project coordinates to EPSG:5070 before computing distances using the
  existing \texttt{redist-map} projection utilities.

\item \textbf{Validate on FL and LA}:
  Florida and Louisiana are the strongest tests because of their complex
  coastal geometries.
  Compare Phase~1 and Phase~2 \cvd{} PP on these states to confirm the
  expected improvement.

\item \textbf{Expose via CLI}:
  \texttt{--cvd-metric geographic} selects Phase~2.
  The default remains \texttt{graph-distance} for backward compatibility.
\end{enumerate}

\subsection{Connection to CVT Theory}

Phase~2 geographic \cvd{} is directly the continuous CVT construction of
Du et al.~\citeyear{du1999} applied to the discrete redistricting domain.
With Euclidean distance, the Phase~2 algorithm satisfies the two CVT conditions:
(1) each tract is assigned to its nearest seed (Voronoi condition);
(2) each seed is at the population-weighted centroid of its district
(centroidal condition, up to the discretisation of approximating a continuous
centroid by the nearest actual tract).

The convergence guarantee of Proposition~\ref{prop:convergence} applies to
Phase~2 without modification: Euclidean Voronoi assignment minimises the sum
of squared distances given fixed seeds, and centroid update minimises the same
sum given fixed assignment.
The discrete rounding (nearest tract to the continuous centroid) introduces
a rounding error, but does not break the monotone energy descent as long as
the error is smaller than the step size.

The convergence rate for Phase~2 is expected to be faster than Phase~1 in
Euclidean-like geographies: the Lloyd energy surface for Euclidean distance
is smoother and has fewer local minima than the graph-distance surface,
because Euclidean distance is a continuous metric while BFS hop count is
discrete and highly non-convex.
